\begin{textsample}{POS Dim 1 – vem\_ed\_16 – Score 13.00 – vem\_ed\_16/t072.txt}  \label{ex:f1_pos_089}
BOAS PRÁTICAS Edilene Gasparini Fernandes, professora de Inglês na Fatec São José do Rio Preto, coordena um Projeto Colaborativo Internacional (PCI / Cesu) desenvolvido com mais seis professores de Inglês, das Fatecs Guaratinguetá, Pompeia, Carapicuíba, Bragança Paulista e Mogi das Cruzes, em parceria com Nitya Sethuraman, \textbf{que} leciona Psicologia da Linguagem na Universidade de Michigan, campus de Dearborn.
Os diferenciais desse PCI, realizado em duas edições (segundo semestre de 2021 e segundo semestre de 2022), \textbf{são} o envolvimento de alunos fluentes em inglês e a execução do projeto exclusivamente por meio de sessões \textbf{síncronas}.
Geralmente, um PCI alterna \textbf{encontros} \textbf{síncronos} — como o quebra-gelo e a apresentação dos trabalhos finais — e atividades assíncronas de \textbf{desenvolvimento} das \textbf{tarefas}."A \textbf{interação} privilegiou alunos fluentes e abordou Psicologia da Linguagem, área de atuação da professora Nitya", relata Edilene.
As etapas do PCI no segundo semestre de 2022 tiveram como temas: 1.
Quebra gelo: Inglês como linguagem internacional; 2.
Linguagem Corporal / Linguagem x Comunicação; 3.
Dialetos, sotaques e identidades; 4.
Bilinguismo / Multilingualismo.
Para desenvolver o debate, formaram-se grupos mistos com quatro brasileiros e seis norte-americanos. \textbf{Todos} os \textbf{encontros} se deram de \textbf{forma} \textbf{síncrona}, entre 3 de outubro e 4 de dezembro de 2022.
Em novembro, a \textbf{diferença} de fuso \textbf{horário}, \textbf{que} \textbf{era} de uma hora, passou a \textbf{ser} de duas horas por causa do \textbf{horário} de inverno nos Estados Unidos."Felizmente, os alunos americanos conseguiram adiantar suas atividades sem alterar as nossas", \textbf{conta} Edilene.
Durante os \textbf{encontros}, os estudantes debatiam o assunto pré-estabelecido e, com questionamentos \textbf{propostos} pela professora Nitya, os docentes brasileiros organizavam as discussões ao \textbf{longo} de 60 minutos.
Segundo Edilene, um desafio de peso \textbf{foi}"comprometer alunos \textbf{que} já têm dispensa na \textbf{disciplina} de Inglês pelo teste de proficiência realizado pelas Fatecs (NEPLE) e se mostram interessados em participar do PCI. \textbf{É} preciso ' ter samba no pé ' para mantê-los até o final".
Isso porque não há vínculos com uma \textbf{disciplina} e, portanto, não há nota para as atividades.
Para motivar a \textbf{turma}, Edilene buscou mostrar a ligação entre os tópicos propostos e a \textbf{realidade} dos estudantes."Por exemplo, o uso dos emojis na comunicação diária ou as crenças \textbf{que} carregamos sobre os sotaques regionais".
Edilene conclui:"Somos um povo muito elogiado dentro dos \textbf{programas} de Intercâmbio Virtual.
O sucesso vem do nosso comprometimento e boa vontade".

% matched lemmas: contar, desenvolvimento, diferença, disciplina, encontro, forma, horário, interação, longo, programa, propor, que, realidade, ser, síncrono, tarefa, todo, turma
\end{textsample}
